\section*{अध्याय \ref{ch:g8:circuit-laws} --- परिपथों में धारा और विभवांतर के नियम}

\begin{solution}{exo:g8:circuit-laws:1}
श्रेणी वाला धारा नियम: एक फंदा, और हर जगह एक ही \omterm{def:g8:intensity-ammeter:intensity}{तीव्रता}। \omterm{def:g5:series-parallel-circuits:parallel}{संधि} का
नियम: भीतर जाती धाराएँ बाहर आती धाराओं के बराबर। श्रेणी वाला
\omterm{def:g8:voltage-voltmeter:voltage}{विभवांतर} नियम: \omterm{def:g3:first-electric-circuit:battery}{बैटरी} का \omterm{def:g8:voltage-voltmeter:voltage}{विभवांतर} फंदे भर के उपकरणों के हिस्सों के जोड़
के बराबर। समांतर वाला \omterm{def:g8:voltage-voltmeter:voltage}{विभवांतर} नियम: उन्हीं दो संधियों पर फैली शाखाएँ
एक ही \omterm{def:g8:voltage-voltmeter:voltage}{विभवांतर} बाँटती हैं।
\end{solution}

\begin{solution}{exo:g8:circuit-laws:2}
$0.8 - 0.55 = \qty{0.25}{A}$.
\end{solution}

\begin{solution}{exo:g8:circuit-laws:3}
$9 - 2.5 - 3.8 = \qty{2.7}{V}$.
\end{solution}

\begin{solution}{exo:g8:circuit-laws:4}
हर एक पर \qty{12}{V} --- यानी समांतर वाला \omterm{def:g8:voltage-voltmeter:voltage}{विभवांतर} नियम। वही नियम घर
की समांतर तारबंदी के हर सॉकेट पर एक जैसी \qty{230}{V} रख देता है।
\end{solution}

\begin{solution}{exo:g8:circuit-laws:5}
जैसे $I_C$: $I_{\text{बैटरी}}$ और $I_B$ पता हो जाने के बाद \omterm{def:g5:series-parallel-circuits:parallel}{संधि} का
नियम $I_C = \qty{0.2}{A}$ तय कर देता है --- वहाँ भेजा गया कोई भी
उपकरण सिर्फ़ पुष्टि ही कर सकता था (इसी तरह $U_C$ भी, जिसे $U_B$ पता
होते ही समांतर वाला \omterm{def:g8:voltage-voltmeter:voltage}{विभवांतर} नियम तय कर देता है)।
\end{solution}

\begin{solution}{exo:g8:circuit-laws:6}
\qty{2.7}{V} वाला लैंप: बराबर धारा पर बड़ा \omterm{def:g8:voltage-voltmeter:voltage}{विभवांतर} हिस्सा ही ज़्यादा
विरोध करने वाले की पहचान है। अगला अध्याय उसी विरोध को मापता है ---
प्रतिरोध।
\end{solution}

\begin{solution}{exo:g8:circuit-laws:7}
$2.2 + 1.9 + 1.4 = \qty{5.5}{V}$, जबकि थैली \qty{6}{V} की है: आधा
\omterm{def:g8:voltage-voltmeter:voltage}{वोल्ट} भटका हुआ --- ऐसे पाठों के लिए यह उपकरण के ईमानदार छितराव से आगे
है; कोई सुई-पट्टी ग़लत पढ़ी गई (या कोई हिस्सा मापा ही नहीं गया)।
\end{solution}

\begin{solution}{exo:g8:circuit-laws:8}
बाहर जाती $0.4 + 0.3 + 0.3 = \qty{1.0}{A}$, जबकि पहुँचती
\qty{0.9}{A}: दोराहा \qty{0.1}{A} ``बना'' रहा है --- \omterm{def:g5:series-parallel-circuits:parallel}{संधि} के नियम ने
दोषी ठहरा दिया।
\end{solution}

\begin{solution}{exo:g8:circuit-laws:9}
फंदे का ख़र्च जुड़कर \omterm{def:g3:first-electric-circuit:battery}{बैटरी} की \qty{4.5}{V} बनना ही चाहिए: खुला \omterm{ex:g3:first-electric-circuit:switch}{स्विच}
\qty{4.5}{V} $+$ ख़ाली लैंप \qty{0}{V} $= 4.5$ --- बराबर। (खटके के बाद
हिस्से जगह बदल लेते हैं: \omterm{ex:g3:first-electric-circuit:switch}{स्विच} क़रीब शून्य पर, लैंप क़रीब \qty{4.5}{V}
पर --- फिर बराबर।)
\end{solution}

\begin{solution}{exo:g8:circuit-laws:10}
पहले: $3 \times 1.5 = \qty{4.5}{V}$. बाद में: $4.5 - 1.5 =
\qty{3.0}{V}$ --- उलटा बैठा सेल घट जाता है। लैंप अपनी \qty{4.5}{V}
वाली दमक से उतरकर \qty{3}{V} वाली दमक पर आ जाता है।
\end{solution}

\begin{solution}{exo:g8:circuit-laws:11}
पहली \omterm{def:g5:series-parallel-circuits:parallel}{शाखा} \omterm{def:g3:first-electric-circuit:battery}{बैटरी} पर फैली है: $U_E + U_F = \qty{6}{V}$ (\omterm{def:g5:series-parallel-circuits:parallel}{शाखा} के
साथ-साथ श्रेणी वाला \omterm{def:g8:voltage-voltmeter:voltage}{विभवांतर} नियम), इसलिए $U_F = 6 - 2.2 =
\qty{3.8}{V}$. दूसरी \omterm{def:g5:series-parallel-circuits:parallel}{शाखा} उन्हीं संधियों पर फैली है: $U_G =
\qty{6}{V}$ (समांतर वाला \omterm{def:g8:voltage-voltmeter:voltage}{विभवांतर} नियम)।
\end{solution}

\begin{solution}{exo:g8:circuit-laws:12}
$I_R = 0.6 - 0.45 = \qty{0.15}{A}$ (\omterm{def:g5:series-parallel-circuits:parallel}{संधि} का नियम)। $U_R = U_Q =
\qty{2.3}{V}$ (समांतर वाला \omterm{def:g8:voltage-voltmeter:voltage}{विभवांतर} नियम)। $U_P = 6 - 2.3 =
\qty{3.7}{V}$ (श्रेणी वाला \omterm{def:g8:voltage-voltmeter:voltage}{विभवांतर} नियम)। दोषी ठहराए जाने लायक एक
फ़ालतू पाठ: ``$I_R = \qty{0.2}{A}$'' --- तब \omterm{def:g5:series-parallel-circuits:parallel}{संधि} का बहीखाता
$0.45 + 0.2 = 0.65 \neq 0.6$ दिखाता; या ``$U_R = \qty{2.6}{V}$'' ---
समांतर शाखाएँ आपस में असहमत हो ही नहीं सकतीं।
\end{solution}

\begin{solution}{pb:g8:circuit-laws:1}
\textbf{1.} श्रेणी। \omterm{def:g4:series-circuits:series}{श्रेणी में} दोनों लैंप उस क़तार का एक के बाद एक
विरोध करते हैं: फंदे की धारा अकेले लैंप की \qty{0.4}{A} से काफ़ी नीचे
गिर जाती है --- और मापी हुई \qty{0.2}{A} इसी में बैठती है। (समांतर तो
दो पूरी शाखाएँ \emph{जोड़} देता।)
\textbf{2.} समांतर: क़रीब \qty{0.4}{A} वाली दो शाखाएँ --- \omterm{def:g3:first-electric-circuit:battery}{बैटरी}
से मोटे तौर पर \qty{0.8}{A}।
\textbf{3.} श्रेणी वाला \omterm{def:g8:voltage-voltmeter:voltage}{विभवांतर} नियम \qty{6}{V} को दो एक जैसे लैंपों
में बाँट देता है: हर एक पर क़रीब \qty{3}{V}।
\textbf{4.} \omterm{def:g5:series-parallel-circuits:parallel}{समांतर में} --- तब हर लैंप पूरी \qty{6}{V} पर फैला होता।
झरोखे से देखने पर श्रेणी वाली जोड़ी मंद दमकती है: हर लैंप अपने आम
धक्के के आधे पर भूखा बैठा है।
\textbf{5.} $I_Z = 0.5 - 0.35 = \qty{0.15}{A}$ (\omterm{def:g5:series-parallel-circuits:parallel}{संधि} का नियम);
$I_X = \qty{0.5}{A}$ --- X बिना बँटी सड़क पर बैठा है (श्रेणी वाला धारा
नियम)।
\textbf{6.} $U_{\text{जोड़ी}} = 6 - 2.8 = \qty{3.2}{V}$ (फंदे वाला
नियम); $U_Y = U_Z = \qty{3.2}{V}$ (समांतर वाला नियम)।
\textbf{7.} वही धक्का, अलग-अलग क़तारें: दोनों लैंप धारा का असमान
विरोध करते हैं --- और Z, जो उसी \omterm{def:g8:voltage-voltmeter:voltage}{विभवांतर} पर कम ढोता है, ज़्यादा विरोध
करने वाला है (यानी ज़्यादा कठिन सड़क)।
\textbf{8.} गिरेगी: Y की \omterm{def:g5:series-parallel-circuits:parallel}{शाखा} जाने से एक सड़क घट जाती है; और बची हुई
Z वाली सड़क उस जोड़ी से ज़्यादा विरोध करती है, इसलिए \omterm{def:g3:first-electric-circuit:battery}{बैटरी} कुल कम धारा
चलाती है।
\textbf{9.} पहली सज़ा, फंदे का नियम: \qty{6}{V} की थैली से
$2.5 + 4.1 = \qty{6.6}{V}$ का ख़र्च। दूसरी सज़ा, \omterm{def:g5:series-parallel-circuits:parallel}{संधि} का नियम: भीतर
\qty{0.3}{A}, बाहर $0.2 + 0.15 = \qty{0.35}{A}$.
\textbf{10.} $6 - 2.5 = \qty{3.5}{V}$.
\textbf{11.} $0.2 + 0.15 = \qty{0.35}{A}$.
\textbf{12.} क्योंकि चारों नियम हर ईमानदार परिपथ को ज़रूरत से
ज़्यादा बाँध देते हैं: हर पाठ को बाक़ी पाठों के जोड़ों से सहमत होना ही
पड़ता है, इसलिए गढ़ा हुआ कोई भी अंक कम से कम एक नियम से झूठ बोलेगा ---
और अकेला अंकगणित उसे दोषी ठहरा देता है।
\textbf{13.} जैसे: ``एक सड़क, एक क़तार --- यानी श्रेणी वाला धारा
नियम। हर दोराहे पर क़तारों का हिसाब बराबर, जितनी भीतर उतनी बाहर ---
यानी \omterm{def:g5:series-parallel-circuits:parallel}{संधि} का नियम। किसी भी फंदे के गिर्द \omterm{def:g3:first-electric-circuit:battery}{बैटरी} का जलप्रपात काम करने
वाले आख़िरी बूँद तक ख़र्च कर देते हैं --- श्रेणी वाला \omterm{def:g8:voltage-voltmeter:voltage}{विभवांतर}
नियम। और उन्हीं दो संधियों के बीच लटकी शाखाएँ एक ही ऊँचाई से पीती हैं
--- यानी समांतर वाला \omterm{def:g8:voltage-voltmeter:voltage}{विभवांतर} नियम।''
\end{solution}
